\documentclass[12pt,a4paper]{article}

\usepackage[utf8]{inputenc}
\usepackage[T1]{fontenc}
\usepackage[brazil]{babel}
\usepackage{geometry}
\usepackage{setspace}
\usepackage{enumitem}
\usepackage{hyperref}

\geometry{margin=2.5cm}
\setstretch{1.2}

\title{\textbf{Regulamento de Uso da Biblioteca de Pendrives de Distribuições Linux}\\
GULECD/UFPR}
\author{}
\date{}

\begin{document}

\maketitle

\section{Apresentação}

A Biblioteca de Pendrives de Distribuições Linux do GULECD/UFPR é uma iniciativa voltada à disseminação do conhecimento,
 à promoção do software livre e ao incentivo à experimentação de sistemas operacionais baseados em Linux.

O projeto consiste na disponibilização de pendrives contendo imagens de instalação de distribuições Linux, preparados de forma padronizada. 
Esses dispositivos podem ser compartilhados entre os participantes do grupo, permitindo
  acesso facilitado a diferentes distribuições sem a necessidade de download individual.

\section{Objetivo}

A biblioteca tem como objetivos principais:

\begin{itemize}
    \item Facilitar o acesso a distribuições Linux para estudantes e membros da comunidade;
    \item Promover a cultura de software livre e colaboração;
    \item Incentivar a experimentação prática em ambientes controlados
\end{itemize}

\section{Funcionamento da Biblioteca}

O funcionamento da biblioteca baseia-se no compartilhamento organizado dos pendrives entre os participantes.

\begin{itemize}
    \item O usuário pode retirar um pendrive na secretaria do curso de Estatística e Ciência de Dados ou diretamente com outro membro do grupo;
    \item O período máximo de permanência com o pendrive é de \textbf{7 dias corridos};
    \item Após esse período, o pendrive deve ser devolvido na secretaria do curso ou repassado a outro membro interessado;
    \item Toda retirada e devolução deve ser registrada por meio do QR code presente no dispositivo.
\end{itemize}

\section{Procedimento de Retirada e Devolução}

Cada pendrive possui um QR code associado a um formulário de controle.

\begin{itemize}
    \item No momento da retirada, o usuário deve escanear o QR code e preencher o formulário correspondente;
    \item No momento da devolução ou repasse, o mesmo procedimento deve ser realizado;
    \item O correto preenchimento é essencial para rastreabilidade e organização da biblioteca.
\end{itemize}

\section{Uso Adequado das Pendrives}

As pendrives da biblioteca foram configuradas como dispositivos de leitura (read-only) com imagens padronizadas.

\begin{itemize}
    \item Não é permitido utilizar a pendrive para armazenamento de arquivos pessoais;
    \item Não é permitido modificar qualquer conteúdo presente na pendrive;
    \item Não é permitido utilizar a pendrive para execução de atividades que comprometam sua integridade;
    \item Não é permitido utilizar a pendrive para quaisquer atividades que infrinjam a legislação brasileira;
    \item A utilização deve se restringir à inicialização (boot), experimentação no ambiente Linux live e instalação das distribuições disponíveis.
\end{itemize}

\section{Verificação de Integridade}

Cada pendrive possui checksums criptográficos previamente gerados e assinados digitalmente.

\begin{itemize}
    \item O usuário deve, sempre que possível, verificar os checksums antes da utilização;
    \item Caso qualquer divergência seja identificada, a pendrive \textbf{não deve ser utilizada};
    \item O problema deve ser comunicado imediatamente à coordenação do grupo;
    \item Ferramentas e instruções para verificação de integridade são disponibilizadas no repositório do projeto em https://github.com/GUL-Estatistica-e-Ciencia-de-Dados-UFPR/biblioteca-distros.
\end{itemize}

\section{Segurança e Responsabilidade do Usuário}

Foram empregados os melhores esforços técnicos para garantir a integridade e segurança das pendrives, incluindo:

\begin{itemize}
    \item uso de imagens oficiais;
    \item configuração para minimizar persistência e alterações;
    \item geração e assinatura criptográfica de checksums;
    \item processos de validação e verificação.
\end{itemize}

Entretanto, \textbf{o usuário é responsável por verificar a integridade da pendrive antes do uso}. Isso inclui:

\begin{itemize}
    \item validar checksums;
    \item verificar assinaturas digitais dos arquivos que contém as checksums;
    \item garantir que o dispositivo não foi adulterado.
\end{itemize}

\section{Isenção de Responsabilidade}

O GULECD/UFPR não se responsabiliza por:

\begin{itemize}
    \item perda de dados;
    \item danos a hardware ou software;
    \item falhas decorrentes do uso das pendrives;
    \item quaisquer consequências diretas ou indiretas da utilização dos dispositivos.
\end{itemize}

O uso das pendrives é de inteira responsabilidade do usuário.

\section{Boas Práticas Recomendadas}

Para uma utilização segura e eficiente, recomenda-se:

\begin{itemize}
    \item realizar backup de dados importantes antes de qualquer instalação;
    \item evitar uso em sistemas críticos sem validação prévia;
    \item manter boas práticas de segurança digital.
\end{itemize}

\section{Conservação dos Dispositivos}

Para garantir a durabilidade das pendrives:

\begin{itemize}
    \item evite impactos físicos;
    \item não exponha a líquidos ou temperaturas extremas;
    \item não remova o dispositivo durante operações de leitura;
    \item mantenha o dispositivo limpo e protegido.
\end{itemize}

\section{Conduta e Colaboração}

A biblioteca depende da colaboração dos participantes.

\begin{itemize}
    \item respeite os prazos de uso;
    \item registre corretamente as movimentações;
    \item reporte problemas técnicos;
    \item contribua com sugestões de melhoria;
    \item respeite os demais usuários do sistema.
\end{itemize}

\section{Disposições Finais}

O uso das pendrives implica na concordância com este regulamento.

O GULECD/UFPR poderá atualizar este documento conforme necessário para aprimorar o funcionamento da biblioteca.

\end{document}